\documentclass[10pt,a4paper]{amsart}
\usepackage{amsmath,amssymb,amsthm,graphicx}
% Retained preamble for the Stacks Project chapter morphisms.tex (Morphisms of Schemes).
% The chapter file pinned-original-source.tex carries only \input{preamble} on its first line, so no
% preamble text lives inside the chapter itself.  The Stacks Project uses one single shared preamble for
% every chapter; the snapshot stages it verbatim in sources/stacks-chow/preamble.tex.  The block below is
% copied mechanically, line for line, from that file's lines 14-18 (the Xy-pic setup) and lines 154-217
% (the native theorem declarations together with the mathematical macros).  Nothing was added, removed or
% redefined.  The wrapper class and the amsmath/amssymb/amsthm/graphicx base packages are supplied by the
% assembler, so the class/\externaldocument part of the shared preamble is not repeated here.
% For commutative diagrams we use Xy-pic
\usepackage[all]{xy}

% We use 2cell for 2-commutative diagrams.
\xyoption{2cell}
% Theorem environments.
%
\theoremstyle{plain}
\newtheorem{theorem}[subsection]{Theorem}
\newtheorem{proposition}[subsection]{Proposition}
\newtheorem{lemma}[subsection]{Lemma}

\theoremstyle{definition}
\newtheorem{definition}[subsection]{Definition}
\newtheorem{example}[subsection]{Example}
\newtheorem{exercise}[subsection]{Exercise}
\newtheorem{situation}[subsection]{Situation}

\theoremstyle{remark}
\newtheorem{remark}[subsection]{Remark}
\newtheorem{remarks}[subsection]{Remarks}

\numberwithin{equation}{subsection}

% Macros
%
\def\lim{\mathop{\mathrm{lim}}\nolimits}
\def\colim{\mathop{\mathrm{colim}}\nolimits}
\def\Spec{\mathop{\mathrm{Spec}}}
\def\Hom{\mathop{\mathrm{Hom}}\nolimits}
\def\Ext{\mathop{\mathrm{Ext}}\nolimits}
\def\SheafHom{\mathop{\mathcal{H}\!\mathit{om}}\nolimits}
\def\SheafExt{\mathop{\mathcal{E}\!\mathit{xt}}\nolimits}
\def\Sch{\mathit{Sch}}
\def\Mor{\mathop{\mathrm{Mor}}\nolimits}
\def\Ob{\mathop{\mathrm{Ob}}\nolimits}
\def\Sh{\mathop{\mathit{Sh}}\nolimits}
\def\NL{\mathop{N\!L}\nolimits}
\def\CH{\mathop{\mathrm{CH}}\nolimits}
\def\proetale{{pro\text{-}\acute{e}tale}}
\def\etale{{\acute{e}tale}}
\def\QCoh{\mathit{QCoh}}
\def\Ker{\mathop{\mathrm{Ker}}}
\def\Im{\mathop{\mathrm{Im}}}
\def\Coker{\mathop{\mathrm{Coker}}}
\def\Coim{\mathop{\mathrm{Coim}}}

% Boxtimes
%
\DeclareMathSymbol{\boxtimes}{\mathbin}{AMSa}{"02}

%
% Macros for moduli stacks/spaces
%
\def\QCohstack{\mathcal{QC}\!\mathit{oh}}
\def\Cohstack{\mathcal{C}\!\mathit{oh}}
\def\Spacesstack{\mathcal{S}\!\mathit{paces}}
\def\Quotfunctor{\mathrm{Quot}}
\def\Hilbfunctor{\mathrm{Hilb}}
\def\Curvesstack{\mathcal{C}\!\mathit{urves}}
\def\Polarizedstack{\mathcal{P}\!\mathit{olarized}}
\def\Complexesstack{\mathcal{C}\!\mathit{omplexes}}
% \Pic is the operator that assigns to X its picard group, usage \Pic(X)
% \Picardstack_{X/B} denotes the Picard stack of X over B
% \Picardfunctor_{X/B} denotes the Picard functor of X over B
\def\Pic{\mathop{\mathrm{Pic}}\nolimits}
\def\Picardstack{\mathcal{P}\!\mathit{ic}}
\def\Picardfunctor{\mathrm{Pic}}
\def\Deformationcategory{\mathcal{D}\!\mathit{ef}}

% --- appended 2026-10-12 (Chief): the project own definitions of the reference/slogan/history
% display environments, copied mechanically from the shared Stacks preamble lines 8-12
% (\usepackage{verbatim} plus the three \newenvironment lines, which the chapter preamble keeps
% verbatim from the project preamble). Nothing added, removed or redefined: these environments
% carry no mathematical content and the project itself typesets them as comments. Without these
% three lines a counted statement whose own text uses them fails with Environment slogan undefined.
\usepackage{verbatim}
\newenvironment{reference}{\comment}{\endcomment}
\newenvironment{slogan}{\comment}{\endcomment}
\newenvironment{history}{\comment}{\endcomment}
\usepackage{cleveref}
\crefname{originallemma}{original lemma}{original lemmas}
\Crefname{originallemma}{Original lemma}{Original lemmas}
\makeatletter
\newlabel{algebra-lemma-integral-closure-in-normal}{{034L}{source}{}{sourcepointer-1}{}}
\newlabel{algebra-lemma-integral-closure-in-normal@cref}{{[originallemma][0][]034L}{[1][0][]source}{}{}{}}
\makeatother
\crefname{originallemma}{original lemma}{original lemmas}
\Crefname{originallemma}{Original lemma}{Original lemmas}
\makeatletter
\newlabel{lemma-integral-closure}{{035F}{source}{}{sourcepointer-2}{}}
\newlabel{lemma-integral-closure@cref}{{[originallemma][0][]035F}{[1][0][]source}{}{}{}}
\makeatother
\crefname{originallemma}{original lemma}{original lemmas}
\Crefname{originallemma}{Original lemma}{Original lemmas}
\makeatletter
\newlabel{lemma-normalization-in-disjoint-union}{{03GO}{source}{}{sourcepointer-3}{}}
\newlabel{lemma-normalization-in-disjoint-union@cref}{{[originallemma][0][]03GO}{[1][0][]source}{}{}{}}
\makeatother
\crefname{originallemma}{original lemma}{original lemmas}
\Crefname{originallemma}{Original lemma}{Original lemmas}
\makeatletter
\newlabel{lemma-normalization-localization}{{035K}{source}{}{sourcepointer-4}{}}
\newlabel{lemma-normalization-localization@cref}{{[originallemma][0][]035K}{[1][0][]source}{}{}{}}
\makeatother
\crefname{originallemma}{original lemma}{original lemmas}
\Crefname{originallemma}{Original lemma}{Original lemmas}
\makeatletter
\newlabel{properties-lemma-locally-normal}{{033J}{source}{}{sourcepointer-5}{}}
\newlabel{properties-lemma-locally-normal@cref}{{[originallemma][0][]033J}{[1][0][]source}{}{}{}}
\makeatother
\crefname{originallemma}{original lemma}{original lemmas}
\Crefname{originallemma}{Original lemma}{Original lemmas}
\makeatletter
\newlabel{properties-lemma-normal-integral-sections}{{0358}{source}{}{sourcepointer-6}{}}
\newlabel{properties-lemma-normal-integral-sections@cref}{{[originallemma][0][]0358}{[1][0][]source}{}{}{}}
\makeatother
\crefname{originallemma}{original lemma}{original lemmas}
\Crefname{originallemma}{Original lemma}{Original lemmas}
\makeatletter
\newlabel{properties-lemma-normal-locally-finite-nr-irreducibles}{{0357}{source}{}{sourcepointer-7}{}}
\newlabel{properties-lemma-normal-locally-finite-nr-irreducibles@cref}{{[originallemma][0][]0357}{[1][0][]source}{}{}{}}
\makeatother
\title{The normalization of a base in a normal source is a disjoint union of integral normal schemes with dominant structure morphism}
\author{The Stacks Project Authors}
\date{stacks-project:morphisms.tex:local-snapshot-20261012; GNU FDL 1.2 or later}
\begin{document}
\maketitle
\noindent\textit{Editorial note.} All delivered mathematical text is the author's own, copied from the pinned Stacks Project chapter file pinned-original-source.tex; no statement or proof step was rewritten, reordered or removed.\par\medskip
\noindent\textit{Editorial note.} Following the multi-claim convention for this source, source\_url is the source-level URL https://stacks.math.columbia.edu/ and the source label lemma-normal-normalization is carried in claim\_key and here rather than in a per-claim tag URL.\par\medskip
\noindent\textit{Editorial note.} The preamble is the chapter's shared Stacks Project preamble, retained as sources/stacks-morphisms/retained-preamble.tex: a line-for-line excerpt of the shared Stacks preamble staged with this snapshot. No macro was invented or redefined, and this card adds no preamble text of its own.\par\medskip
\noindent\textit{Editorial note.} Statement span 14169-14180 and proof span 14182-14218 come from the halfA pool card operations/pool-cards/stacks-morphisms-halfA.json and were re-read in the pinned chapter file before this card was written.\par\medskip
\noindent\textit{Editorial note.} The proof twice refers to a topological argument that the author leaves out ('which we omit' and 'omitted'), once for the local finiteness of irreducible components of the normalization and once for the bijection on irreducible components; both phrases are kept verbatim, and in each case the geometric input (the description of the normalization over affine opens) is established in the counted span.\par\medskip
\noindent\textit{Editorial note.} 7 cross-references occur inside the counted spans; they are listed in external\_references as navigation pointers only.\par\medskip
\noindent\textit{Allowed standard prerequisites (editorial).}\begin{itemize}
\item The compatibility of normalization with localization, cited from the same chapter. Cited at proof lines 14183-14185. Reduces the statement over an affine open of the base to the normalization of that ring in the preimage.
\item A normal scheme whose quasi-compact opens have finitely many irreducible components is a disjoint union of normal integral schemes, cited from the properties chapter. Cited at proof lines 14189-14192 and 14208-14211. Decomposes the preimage of an affine open into finitely many normal integral pieces and gives the global conclusion.
\item The compatibility of normalization with disjoint unions, established by a previous card of this lane, and the description of the normalization by integral closures. Cited at proof lines 14193-14196 and 14200-14203. Splits the normalization over the affine open into the normalizations in the individual pieces.
\item The ring of global functions of a normal integral affine scheme is a normal domain, and the integral closure of a ring in a normal domain is a normal domain. Cited at proof lines 14197-14206. Shows that each piece of the normalization is again a normal integral scheme.
\end{itemize}\medskip
\ifdefined\hypertarget\hypertarget{sourcepointer-1}{}\fi
\noindent\textit{Source reference.} Original lemma 034L: focused-seeds/AGP/sources/stacks-snapshot-20261011/tags.txt, source lines 4096--4096. Same-chapter reference used in the retained author text at source line(s) 14200. It is supplied only as a navigation pointer so that the retained text compiles; the counted item delivers the author's own statement and the author's complete proof, and no part of the referenced item is copied in or substituted for the counted argument.\par\medskip
\ifdefined\hypertarget\hypertarget{sourcepointer-2}{}\fi
\noindent\textit{Source reference.} Original lemma 035F: focused-seeds/AGP/sources/stacks-snapshot-20261011/tags.txt, source lines 4126--4126. Same-chapter reference used in the retained author text at source line(s) 13884, 13969, 14199, 14489. It is supplied only as a navigation pointer so that the retained text compiles; the counted item delivers the author's own statement and the author's complete proof, and no part of the referenced item is copied in or substituted for the counted argument.\par\medskip
\ifdefined\hypertarget\hypertarget{sourcepointer-3}{}\fi
\noindent\textit{Source reference.} Original lemma 03GO: focused-seeds/AGP/sources/stacks-snapshot-20261011/tags.txt, source lines 4535--4535. Same-chapter reference used in the retained author text at source line(s) 14192, 14455. It is supplied only as a navigation pointer so that the retained text compiles; the counted item delivers the author's own statement and the author's complete proof, and no part of the referenced item is copied in or substituted for the counted argument.\par\medskip
\ifdefined\hypertarget\hypertarget{sourcepointer-4}{}\fi
\noindent\textit{Source reference.} Original lemma 035K: focused-seeds/AGP/sources/stacks-snapshot-20261011/tags.txt, source lines 4131--4131. Same-chapter reference used in the retained author text at source line(s) 14087, 14185, 14659. It is supplied only as a navigation pointer so that the retained text compiles; the counted item delivers the author's own statement and the author's complete proof, and no part of the referenced item is copied in or substituted for the counted argument.\par\medskip
\ifdefined\hypertarget\hypertarget{sourcepointer-5}{}\fi
\noindent\textit{Source reference.} Original lemma 033J: focused-seeds/AGP/sources/stacks-snapshot-20261011/tags.txt, source lines 4058--4058. Same-chapter reference used in the retained author text at source line(s) 14207. It is supplied only as a navigation pointer so that the retained text compiles; the counted item delivers the author's own statement and the author's complete proof, and no part of the referenced item is copied in or substituted for the counted argument.\par\medskip
\ifdefined\hypertarget\hypertarget{sourcepointer-6}{}\fi
\noindent\textit{Source reference.} Original lemma 0358: focused-seeds/AGP/sources/stacks-snapshot-20261011/tags.txt, source lines 4119--4119. Same-chapter reference used in the retained author text at source line(s) 14195. It is supplied only as a navigation pointer so that the retained text compiles; the counted item delivers the author's own statement and the author's complete proof, and no part of the referenced item is copied in or substituted for the counted argument.\par\medskip
\ifdefined\hypertarget\hypertarget{sourcepointer-7}{}\fi
\noindent\textit{Source reference.} Original lemma 0357: focused-seeds/AGP/sources/stacks-snapshot-20261011/tags.txt, source lines 4118--4118. Same-chapter reference used in the retained author text at source line(s) 14191, 14212, 14660. It is supplied only as a navigation pointer so that the retained text compiles; the counted item delivers the author's own statement and the author's complete proof, and no part of the referenced item is copied in or substituted for the counted argument.\par\medskip
\begin{lemma}
\label{lemma-normal-normalization}
Let $f : Y \to X$ be a quasi-compact and quasi-separated morphism
of schemes. Let $X'$ be the normalization of $X$ in $Y$. Assume
\begin{enumerate}
\item $Y$ is a normal scheme,
\item quasi-compact opens of $Y$ have finitely many irreducible components.
\end{enumerate}
Then $X'$ is a disjoint union of integral normal schemes. Moreover,
the morphism $Y \to X'$ is dominant and induces a bijection of
irreducible components.
\end{lemma}

\begin{proof}
Let $U \subset X$ be an affine open. Consider the
inverse image $U'$ of $U$ in $X'$.
Set $V = f^{-1}(U)$. By Lemma \ref{lemma-normalization-localization}
we $V \to U' \to U$ is the normalization of $U$ in $V$. Say
$U = \Spec(A)$. Then $V$ is quasi-compact, and hence has a finite number of
irreducible components by assumption. Hence
$V = \coprod_{i = 1, \ldots n} V_i$ is a finite disjoint union of
normal integral schemes by
Properties, Lemma \ref{properties-lemma-normal-locally-finite-nr-irreducibles}.
By Lemma \ref{lemma-normalization-in-disjoint-union}
we see that $U' = \coprod_{i = 1, \ldots, n} U_i'$,
where $U'_i$ is the normalization of $U$ in $V_i$.
By Properties, Lemma \ref{properties-lemma-normal-integral-sections}
we see that $B_i = \Gamma(V_i, \mathcal{O}_{V_i})$ is a normal domain.
Note that $U_i' = \Spec(A_i')$, where $A_i' \subset B_i$
is the integral closure of $A$ in $B_i$, see
Lemma \ref{lemma-integral-closure}. By
Algebra, Lemma \ref{algebra-lemma-integral-closure-in-normal}
we see that $A_i' \subset B_i$ is a normal domain.
Hence $U' = \coprod U_i'$ is a finite union of normal integral schemes
and hence is normal.

\medskip\noindent
As $X'$ has an open covering by the schemes $U'$ we conclude from
Properties, Lemma \ref{properties-lemma-locally-normal} that $X'$ is normal.
On the other hand, each $U'$ is a finite disjoint union of irreducible
schemes, hence every quasi-compact open of $X'$ has finitely many irreducible
components (by a topological argument which we omit). Thus $X'$
is a disjoint union of normal integral schemes by
Properties, Lemma \ref{properties-lemma-normal-locally-finite-nr-irreducibles}.
It is clear from the description of $X'$ above that $Y \to X'$
is dominant and induces a bijection on irreducible components
$V \to U'$ for every affine open $U \subset X$. The bijection of
irreducible components for the morphism $Y \to X'$
follows from this by a topological argument (omitted).
\end{proof}

\end{document}
